\documentclass[12pt,a4paper,twoside]{article}
\usepackage[english]{babel}
\usepackage{amsmath}
\usepackage{amssymb}
\usepackage{amsfonts}
\usepackage{graphicx}
\usepackage[utf8]{inputenc}
\usepackage[hidelinks, final]{hyperref}
\usepackage[printonlyused]{acronym}
\usepackage{color}
\usepackage{transparent}
\graphicspath{{img/}}
\usepackage[%
backend=biber,
url=false,
style=numeric,
sorting = none,
maxnames=4,
minnames=3,
maxbibnames=99,
giveninits,
uniquename=init]{biblatex}
\addbibresource{bachelor_thesis_julien_caselmann.bib}

\setlength{\voffset}{-28.4mm}
\setlength{\hoffset}{-1in}
\setlength{\topmargin}{20mm}
\setlength{\oddsidemargin}{25mm}
\setlength{\evensidemargin}{25mm}
\setlength{\textwidth}{160mm}

\setlength{\parindent}{0pt}

\setlength{\textheight}{235mm}
\setlength{\footskip}{20mm}
\setlength{\headsep}{50pt}
\setlength{\headheight}{0pt}


\begin{document}
\pagestyle{empty}
%%%% Title page
\begin{titlepage}
\begin{center}
\includegraphics{img/TUMlblack.png}\\[3mm]
\sf
{\Large
  Technische Universit\"at M\"unchen\\[5mm]
  Department of Mathematics\\[8mm]
}
\normalsize
\includegraphics{img/TUMlMblack.png}\\[15mm]

Bachelor's Thesis\\[15mm]

{\Huge
  The Dynamics of Vaccination Hesitancy
}
\bigskip

\normalsize

Julien Caselmann
\end{center}
\vspace*{75mm}

Supervisor: Prof. Dr. Johannes M\"uller
\medskip

Advisor: Prof. Dr. Johannes M\"uller
\medskip

Submission Date: ...

\end{titlepage}
%%%% The following has to be signed by hand!

\vspace*{150mm}

I assure the single handed composition of this bachelor's thesis only supported by declared resources.
\bigskip

Garching,
\newpage
%%%% Summary in german
\section*{Zusammenfassung}
Bei einer in englischer Sprache verfassten Arbeit muss eine Zusammenfassung in deutscher Sprache vorangestellt werden.
Daf\"ur ist hier Platz.

\newpage
\tableofcontents
\newpage

%%%% Page numbering restarts here
\pagenumbering{arabic}
\pagestyle{headings}

\newpage
%if you do not use this, delete following packages in preamble:

%acronym
\section*{Acronyms}
\begin{acronym}[KPMG]%write longest acronym here for nice formatting purposes
	%to use them write \ac{KPMG} or \acp{KPMG for plural in the text}
	\acro{KPMG}{KPMGeneral}
	\acro{SIR}{Susceptibles, Infecteds, Recovereds}
	%if one acronym's plural form is more than just an "s" attached, you have to specify that:
	\acro{SQ}{sexy query}
	\acroplural{SQ}[SQs]{sexy queries}
\end{acronym}

\section{Introduction}
%How to input a inkscape image:
%\begin{figure}
%	\centering %(or not)
%	\def\svgwidth{175pt}
%	\input{foo/bar/file.pdf_tex}
%	\caption{example}
%	\label{fig:example}
%\end{figure}

\section{Mathematical Tools}

\section{Hesitancy Model}
%TODO: write a little intro for this paragraph
Significant work on this topic has been done by Bauch and Bhattacharyya in \cite{Bauch2012}. They presented a behaviour-incidence model with the aim to describe the dynamics of the number of vaccinators $x$ depending on different factors such as vaccine risk, vaccine efficacy and the number of infectious persons. The model reads as follows:
\begin{align}\label{model:bauch}
	\begin{split}
	\frac{dS}{dt} &= \mu\left(1-\epsilon x\right) - \mu S - \beta SI - \tau S\\
	\frac{dI}{dt} &= -\mu I + \beta SI - \gamma I + \tau S\\
	\frac{dx}{dt} &= \kappa x\left(1 - x\right)\left(-\omega + I\right)
	\end{split}
\end{align}
The variables $S$ and $I$ follow the convention from the well-known and well-studied \ac{SIR}-model from \textbf{findsource} %TODO: find a source
and $x$ represents the relative number of vaccinators in the observed population. A lot of different variables occur in \eqref{model:bauch}, so they will be listed here to get a better overview:
\begin{align}\label{table:bauch_model_variables}
	\begin{split}
	\mu &: \textnormal{birth and death rate of the population}\\
	\epsilon &: \textnormal{vaccine efficacy}\\
	\beta &: \textnormal{infection rate}\\
	\tau &: \textnormal{case importation rate}\\
	\gamma &: \textnormal{recovery rate}\\
	\kappa &: \textnormal{scale factor}\\
	\omega &: \textnormal{vaccine penalty}
	\end{split}
\end{align}
The vaccine penalty $\omega$ basically describes the amount of risk a vaccination brings.

The modeling idea is that every person starts in the group of susceptibles $S$ and either stays in $S$, moves to the group of infecteds $I$, or gets the vaccine and immediately jumps into the group of recovereds $R$. As the main interest of this model is the dynamics of the vaccinators $x$ and once an individual gets into $R$, a return to $S$ or $I$ is impossible, it is assumed that the recovered have left the model, thus are not mentioned in \eqref{model:bauch}. The dynamics of $S$ are pretty straight-forward: the only way to get into $S$, is to get born into the population, so the only positive term is the birth rate $\mu$. On the other hand, there are multiple ways to get ot of $S$: an individual might either die $\left(-\mu S\right)$, get infected by one individual of $I$ $\left(-\beta SI\right)$, be born as a child of vaccinators and get vaccined successfully at birth $\left(-\mu\epsilon x\right)$ or come from another population as an infected and immediately jump from $S$ to $I$ $\left(-\tau S\right)$. 

$I$ is designed in a similar way: newly infecteds were either infected in the observed population $\left(+\beta SI\right)$ or came from another population and brougth the disease with them $\left(+\tau S\right)$. Again, an individual leaves $I$ at death $\left(-\mu I\right)$ or after recovering from the disease $\left(-\gamma I\right)$. 

Finally, let us take a look at the vaccinator dynamics%TODO: what exactly does the scale factor kappa do? what would be an example of it?
. They consist of two factors. The first one multiplies the relative number of vaccinators $x$ with that of non-vaccinators $\left(1-x\right)\dots$%TODO: finish this idea
? The second one is the interesting one: it computes the difference between $I$ and the risk $\omega$ of taking the vaccine. It basically depicts the danger emerging from the disease, perceived by the population. This is also called the prevalence of the disease \textbf{maybe source}%TODO: find source
. When there is a high amount of infecteds and the vaccine penalty is low, the difference will be positive and quite big, leading to a higher growth rate of $x$ (more people will become vaccinators and vaccine their children at birth) and vice versa. When $\omega$ and $I$ are approximatively the same, the growth/decrease of $x$ will turn out quite insignificant.

One may notice that the dynamics of the vaccinators $x$ as presented in the model \eqref{model:bauch} are quite basic and not realistic at all. They imply that the only influencing factor in the switch between vaccinator and non-vaccinator is this kind of ``payoff comparison" described by the difference $I - \omega$ and attribute it an extreme power. If for instance,we considered a population with $N$ individuals, almost no vaccinators at a time $t$, say $x(t) \leq 0.05$, a massive amount of infecteds $I \sim N$ and a somewhat decent vaccine risk. Even if the risk was still high, the enormous $I$ would lead to the ``payoff" being quite big. According to this model, a tremendous amount of anti-vaccinators would suddenly become vaccinators, despite their ideologies and fears. Experience tells us that even in such dramatic context, the way from being an extreme ennemy of vaccines to taking vaccination in consideration can be very long %TODO: nice word for very long
\textbf{medical source}%TODO: find source from a huge disease were some freaks didnt vaccine and all died
. Additionally, exchange between those two camps or differently weighted opinions are completely excluded from this model, yet highly present in the real world.

\section{Filter bubbles and modeling approaches to them}
\section{My Model}
\section{Discussion}

\newpage
\printbibliography
\end{document}



